\subsection{Sensitivity to Boarding and Diagnostic Intensity}
\label{sec:res:sens}

Table~\aref{tab:sens} reports, for the \NMediumTypes{} medium-load ED types, whether the best policy of the core design remains the best when boarding is longer or diagnostic intensity changes (capacities unchanged), and how much is lost if the core best policy is kept. With light diagnostics the core best policy loses at most \SensLightDiagnosticsRatpRegretMax\% in RATP and \SensLightDiagnosticsLosAllMeanRegretMax\% in mean LOS. With high boarding the best policy for RATP changes in some ED types, but keeping the core best policy costs on average \SensHighBoardingRatpRegret\% in RATP (at most \SensHighBoardingRatpRegretMax\%) and \SensHighBoardingLosAllMeanRegret\% in mean LOS (at most \SensHighBoardingLosAllMeanRegretMax\%); the best queue rules are then \SensHighBoardingBestQueues. Long boarding hits split flow hardest, because admitted patients from the vertical area board in Main beds: in \SensHighBoardingNearSat{} ED types (\SensHighBoardingNearSatList) the median Main-space utilization across policies reached 0.95 or more.

Heavy diagnostics is the most demanding variant. In \SensHeavyDiagnosticsNearSat{} ED types (\SensHeavyDiagnosticsNearSatList) the median Main-space utilization across policies was 0.95 or more (up to \SensHeavyDiagnosticsUtilMax), because capacities were sized for typical diagnostic use. Close to saturation the core best policy lost up to \SensHeavyDiagnosticsLosAllMeanRegretMax\% in mean LOS; elsewhere it held up, which suggests that the results are least reliable for an ED operating close to the capacity of its Main area.

